\documentclass[12pt]{article}

% Packages
\usepackage[margin=1in]{geometry}
\usepackage{titlesec}
\usepackage{enumitem}
\usepackage{setspace}
\usepackage{helvet}
\usepackage{xcolor}
\usepackage{tcolorbox}
\usepackage{booktabs}
\usepackage{array}
\usepackage{caption}
\usepackage{lmodern}
\usepackage{fancyhdr} % <-- added for footer
\usepackage{graphicx}
\usepackage{subcaption}
% \usepackage{verbatim}

% Caption settings
\captionsetup[table]{labelformat=empty}

% Fonts and spacing
\renewcommand{\familydefault}{\sfdefault}
\setlength{\parskip}{0.8em}
\setlength{\parindent}{0pt}

% Section formatting
\titleformat{\section}
  {\Large\bfseries\color{blue!70!black}}
  {}{0pt}{}

% Custom colors
\definecolor{primary}{HTML}{3A7CA5}
\definecolor{lightgray}{gray}{0.97}

% Footer setup
\pagestyle{fancy}
\fancyhf{} % clear all headers and footers
\renewcommand{\headrulewidth}{0pt}
\renewcommand{\footrulewidth}{0pt}
\fancyfoot[C]{\textit{CS 216 Assignment Report}}

\begin{document}

\begin{titlepage}
\centering
\vspace*{\fill}

\includegraphics[width=0.4\textwidth]{logo3.png} \\[3em]

{\Huge \textbf{Report of Bitcoin Transaction Lab Assignment}} \\
\vfill
\end{titlepage}

\begin{center}
    {\Huge \bfseries Team - NodeSTAR} \\[0.5em]
    \vspace{1em}
    
    \rule{\textwidth}{0.5pt}
\end{center}

\vspace{1.5em}

% Team Members
\begin{tcolorbox}[colback=lightgray, colframe=primary!50, boxrule=0.5pt, arc=3pt,
left=8pt, right=8pt, top=6pt, bottom=6pt]

\textbf{1) \textcolor{blue}{Rudra Chitkara}} \\
\textbf{Roll Number:} 240041031 \\
\textbf{Github:} \texttt{https://github.com/Rudra-Codes}

\vspace{0.8em}

\textbf{2) \textcolor{blue}{Aayush}} \\
\textbf{Roll Number:} 240041001 \\
\textbf{Github:} \texttt{https://github.com/sunno-bhaato}

\vspace{0.8em}

\textbf{3) \textcolor{blue}{Tanish Yadav}} \\
\textbf{Roll Number:} 240041036 \\
\textbf{Github:} \texttt{https://github.com/tanpsi}

\vspace{0.8em}

\textbf{4) \textcolor{blue}{Siddh Nema}} \\
\textbf{Roll Number:} 240002070 \\
\textbf{Github:} \texttt{https://github.com/SiddhaNema}

\end{tcolorbox}

\newpage

% Events Section
\section*{Part 1: Legacy Address Transactions (P2PKH)}

\subsection*{Workflow}
The transaction flow executes in two main stages based on the \texttt{transactions.py} code logic using the `list\_unspent`, `create\_raw\_tx`, `sign\_raw\_tx`, and `broadcast\_tx` commands which map to RPC routines listunspent, createrawtransaction, signrawtransactionwithwallet, sendrawtransaction:

\begin{enumerate}
    \item \textbf{Setup Phase:} We generate three base58-encoded legacy addresses: \texttt{A}, \texttt{B}, and \texttt{C}. Address \texttt{A} is then funded with \textbf{10 BTC} (as defined by \texttt{DEFAULT\_FUND\_A} in \texttt{config.py}) via the \texttt{sendtoaddress} RPC.
    
    \item \textbf{Address A to B:} A raw transaction is created to send \textbf{9 BTC} (from \texttt{DEFAULT\_TX\_A\_TO\_B}) from \texttt{A} to \texttt{B}. We query \texttt{listunspent} for \texttt{A}'s addresses, extract UTXO(s) with at least 9 BTC, calculate the fee (default rate:10 satoshis/vbyte) using place holder fee and feedback loop and build inputs and outputs (9 BTC to B, the rest minus fees as change returned to A). The transaction is signed with A's private key (using the wallet) and broadcasted to the network. We then mine 1 block to confirm it. We also decode the signed transaction using decoderawtransaction which shows the scriptPubKey.
    
    \item \textbf{Address B to C:} Another transaction sends \textbf{8 BTC} from \texttt{B} to \texttt{C}. We run \texttt{listunspent} for \texttt{B} to obtain the UTXO generated by the previous \texttt{A} $\rightarrow$ \texttt{B} transaction. This UTXO (bearing the exact \texttt{txid} from the last step) becomes the input for this new transaction. Similar fee calculations occur, we sign the transaction with B's private key, and broadcast it to the network. We also decode the signed transaction using decoderawtransaction.
\end{enumerate}

\subsection*{Decoded Scripts}
The decoded scripts representing the transaction details are provided below:

\vspace{1em}
\textbf{Decoded Script for A $\rightarrow$ B Transaction:}
\begin{tcolorbox}[colback=lightgray, colframe=primary!50, boxrule=0.5pt, arc=3pt]
    \textit{\includegraphics[width=1\textwidth]{legacy-AtoB.png}}
\end{tcolorbox}

\vspace{1em}
\textbf{Decoded Script for B $\rightarrow$ C Transaction:}
\begin{tcolorbox}[colback=lightgray, colframe=primary!50, boxrule=0.5pt, arc=3pt]
    \textit{\includegraphics[width=1\textwidth]{legacy-BtoC.png}}
\end{tcolorbox}

\subsection*{Challenge and Response Scripts Structure}
The validation of a Legacy P2PKH transaction involves combining the unlocking script (\texttt{scriptSig}) and the locking script (\texttt{scriptPubKey}):
\begin{itemize}
    \item \textbf{Challenge Script (scriptPubKey):} \texttt{OP\_DUP OP\_HASH160 <PubkeyHash> OP\_EQUALVERIFY OP\_CHECKSIG}
    \item \textbf{Response Script (scriptSig):} \texttt{<Signature> <PublicKey>}
\end{itemize}
\textbf{Validation Process:} To spend a P2PKH UTXO, the spender must provide the full public key (which hashes to the \texttt{<PubkeyHash>}) and a valid signature. The network essentially prepends the \texttt{scriptSig} to the \texttt{scriptPubKey}. During validation: 
\begin{enumerate}
    \item \texttt{<Signature>} and \texttt{<PublicKey>} are pushed onto the stack.
    \item \texttt{OP\_DUP} duplicates the \texttt{PublicKey}.
    \item \texttt{OP\_HASH160} hashes the duplicated public key.
    \item \texttt{<PubkeyHash>} is pushed, and \texttt{OP\_EQUALVERIFY} checks if the generated hash equals the provided one. They match, popping both from the stack.
    \item Finally, \texttt{OP\_CHECKSIG} verifies the original signature against the original \texttt{PublicKey} and the transaction data. It evaluates to \texttt{True}, making the spend valid.
\end{enumerate}

\subsection*{Screenshots Evidence for Btcdeb showing script execution in action}
\includegraphics[width=1\textwidth]{btcdeb-legacy-1.png}\\[2em]
\includegraphics[width=1\textwidth]{btcdeb-legacy-2.png}

% \newpage
\begin{center}
    % \includegraphics[width=0.8\textwidth]{placeholder1.png} \\
    \vspace{7cm}
\end{center}

\newpage
\section*{Part 2: P2SH-SegWit Address Transactions}

\subsection*{Workflow}
The transaction flow follows a process identical to Part 1, modified for base58 P2SH address wrappers containing nested SegWit (P2WPKH) programs.
\begin{enumerate}
    \item \textbf{Setup Phase:} Generate three P2SH-SegWit addresses (\texttt{A'}, \texttt{B'}, and \texttt{C'}), starting with "2" (on regtest). Address \texttt{A'} is funded with 10 BTC.
    \item \textbf{Address A' to B':} We create a raw transaction to send 9 BTC to Address \texttt{B'} using the UTXOs for \texttt{A'}. We query listunspent for A’s addresses, extract UTXO(s) with at least 9
BTC, calculate the fee (default rate:10 satoshis/vbyte) using place holder fee and feedback
loop and build inputs and outputs (9 BTC to B, the rest minus fees as change returned to
A). The transaction is signed with A’s private key (using the wallet) and broadcasted to
the network. We then mine 1 block to confirm it. We also decode the signed transaction using decoderawtransaction which shows the scriptPubKey.
    \item \textbf{Address B' to C':} The transaction provides a UTXO locked to \texttt{B'}. We query the UTXOs for \texttt{B'}, isolate the exact \texttt{txid}, and create the next transaction. Upon signing, Bitcoin Core automatically places the SegWit signatures not in the \texttt{scriptSig}, but in the segregated \texttt{witness} data container.
\end{enumerate}

\subsection*{Decoded Scripts}
\vspace{1em}
\textbf{Decoded Script for A' $\rightarrow$ B' Transaction:}
\begin{tcolorbox}[colback=lightgray, colframe=primary!50, boxrule=0.5pt, arc=3pt]
    \textit{\includegraphics[width=1\textwidth]{p2sh-segwit-AtoB.png}}
\end{tcolorbox}

\vspace{1em}
\textbf{Decoded Script for B' $\rightarrow$ C' Transaction:}
\begin{tcolorbox}[colback=lightgray, colframe=primary!50, boxrule=0.5pt, arc=3pt]
    \textit{\includegraphics[width=1\textwidth]{p2sh-segwit-BtoC.png}}
\end{tcolorbox}

\subsection*{Challenge and Response Scripts Structure}
For nested P2SH-P2WPKH transactions, SegWit details are nested within a P2SH construct to maintain backward compatibility with older nodes:
\begin{itemize}
    \item \textbf{Challenge Script (scriptPubKey):} \texttt{OP\_HASH160 <Hash160(RedeemScript)> OP\_EQUAL}
    \item \textbf{Response Script (scriptSig):} \texttt{<0> <20-byte-PubkeyHash>}
    \item \textbf{Witness Stack:} \texttt{<Signature> <PublicKey>}
\end{itemize}
\textbf{Validation Process:} 
\begin{enumerate}
    \item Old nodes view this as a standard P2SH check. They hash the provided \texttt{RedeemScript} (\texttt{0 <20-byte-PubkeyHash>}) from the \texttt{scriptSig} and ensure it matches the 20-byte hash in the \texttt{scriptPubKey}. They stop here, assuming it's valid.
    \item Upgraded SegWit-aware nodes see that the supplied redeem script follows the very specific pattern of \texttt{0} followed by a 20-byte hash (the witness program).
    \item This triggers the secondary validation step: retrieving the external \texttt{witness} stack. The execution context switches to evaluate it exactly like a native SegWit P2WPKH transaction, duplicating the public key, checking its hash against the 20-byte hash in the witness program, and evaluating \texttt{OP\_CHECKSIG} using the provided signature! 
\end{enumerate}

\subsection*{Screenshots Evidence for Btcdeb showing script execution in action}
\includegraphics[width=1\textwidth]{btcdeb-p2sh-segwit-1.png}\\[2em]
\includegraphics[width=1\textwidth]{btcdeb-p2sh-segwit-2.png}
\begin{center}
    % \includegraphics[width=0.8\textwidth]{placeholder2.png} \\
    \vspace{7cm}
\end{center}

\newpage
\section*{Part 3: Analysis and Explanation}

\subsection*{Size Comparison of P2PKH and P2SH-SegWit Transactions}
By observing the differences between Part 1 and Part 2 decoded JSON objects (specifically the \texttt{size}, \texttt{vsize}, and \texttt{weight} fields), we highlight distinct space efficiencies:
\begin{itemize}
    \item \textbf{Legacy P2PKH Size:} The signature script is embedded directly into the transaction input (in the \texttt{scriptSig}), weighing identically as the rest of the transaction. A typical 1-input 1-output P2PKH transaction is heavily penalized. It has a raw size and virtual size mapping of $1:1$ (as shown in image $\sim 225$ bytes translates to $\sim 225$ vBytes and $900$ weight units).
    \item \textbf{P2SH-SegWit Size:} Witness data is segregated from the standard transaction tree structure. While P2SH-SegWit often produces a slightly higher absolute raw byte count because of the P2SH wrapper overhead (e.g., $\sim 247$ bytes), the \textit{virtual size} (vBytes) and \textit{weight} are significantly lower (e.g., $\sim 166$ vBytes and $\sim 661$ weight units).
\end{itemize}
\begin{figure}[h]
    \centering
    
    \begin{subfigure}{0.45\textwidth}
        \centering
        \includegraphics[width=\linewidth]{size_legacy.png}
        \caption{Legacy Size Info}
        \label{fig:first}
    \end{subfigure}
    \hfill
    \begin{subfigure}{0.45\textwidth}
        \centering
        \includegraphics[width=\linewidth]{size_segwit.png}
        \caption{Segwit Size Info}
        \label{fig:second}
    \end{subfigure}

    \caption{Compare how virtual size and fees decreases in segwit as compared to legacy for same transactions size and fees}
    \label{fig:both}
\end{figure}



\subsection*{Benefits of SegWit}
The SegWit (Segregated Witness BIP 141) architecture introduces wide-reaching benefits to the Bitcoin ecosystem:
\begin{enumerate}
    \item \textbf{Lower Transaction Fees:} Users pay fees per virtual byte (vByte). Because witness data (signatures) are given a $75\%$ discount (1 byte = 1 weight unit = 0.25 vBytes), SegWit transactions inherently incur drastically cheaper miner fees to occupy block space.
    \item \textbf{Effective Block Size Increase:} This weight discounting model acts as a soft fork block size bump. A block filled exclusively with SegWit transactions could approximately hold 1.61x more data compared to legacy.
    \item \textbf{Transaction Malleability Fix:} Because the signature data determines its own independent witness hash (`wtxid`) rather than the core transaction ID (`txid`), a third party altering a signature cannot mutate the `txid`. This solidifies unconfirmed transactions, enabling reliable Layer-2 solutions like the Lightning Network.
\end{enumerate}

\end{document}
